\documentclass[10pt,a4paper]{article}
\usepackage[utf8]{inputenc}
\usepackage[spanish]{babel}
\usepackage{amsmath}
\usepackage{amsfonts}
\usepackage{amssymb}
\usepackage{graphicx}
\usepackage[left=2cm,right=2cm,top=2cm,bottom=2cm]{geometry}

\title{\textbf{Dinámica de sistemas físicos}\\Ejercicios: Capítulo 1}
\author{N. Pérez Medina.}
\date{\today}

\begin{document}
\maketitle
\newpage
\section*{Ejercicio 1}
Decide si cada una de las siguientes funciones es \textbf{inyectiva} (\textit{one-to-one}), 
\textbf{sobreyectiva} (\textit{onto}), \textbf{homeomorfismo} o \textbf{difeomorfismo} 
en su \textbf{dominio de definición}:

\begin{enumerate}
    \item $f(x) = x^{5/3}$
    \newpage
    \item $f(x) = x^{4/3}$
    \newpage
    \item $f(x) = 3x + 5$
    \newpage
    \item $f(x) = e^{x}$
    \newpage
    \item $f(x) = \dfrac{1}{x}$
    \newpage
    \item $f(x) = \dfrac{1}{x^{2}}$
\end{enumerate}
\newpage

\section*{Ejercicio 4}
El objetivo de los siguientes ejercicios es construir funciones especiales que serán útiles
más adelante cuando \emph{perturbemos} o modifiquemos ligeramente una función dada.
A estas funciones se les llama \textit{funciones “bump”}.

Defina
\[
B(x)=
\begin{cases}
\exp\!\left(-\dfrac{1}{x^{2}}\right), & \text{si } x>0,\\[6pt]
0, & \text{si } x\le 0.
\end{cases}
\]

\begin{enumerate}
	\item Traza (esboza) la gráfica de $B(x)$.
\end{enumerate}
\newpage

\section*{Ejercicio 7}
Modifica $C(x)$ para construir una \emph{función bump} $C^\infty$ $D(x)$ en el intervalo $[a,b]$, es decir, $D(x)$ satisface:

\begin{enumerate}
    \item $D(x)=1$ para $a \le x \le b$.
    \item $D(x)=0$ para $x<\alpha$ y $x>\beta$, donde $\alpha<a$ y $\beta>b$.
    \item $D'(x)\neq 0$ en los intervalos $(\alpha,a)$ y $(b,\beta)$.
\end{enumerate}



\end{document}